\documentclass[10pt,a4paper]{amsart}
\usepackage[margin=19mm]{geometry}
\usepackage{amsmath,amssymb,mathtools}
\usepackage[scr=rsfs,cal=euler]{mathalfa}
\usepackage{xfrac}
\usepackage[unicode,hidelinks,bookmarks=false]{hyperref}
\newcommand{\ie}{i.e.\ }
\newcommand{\cf}{cf.\ }
\newcommand{\resp}{respectively\ }
\newcommand{\Subsection}{\subsection}
\providecommand{\nobreakdash}{\nobreak\hbox{-}}
\setlength{\emergencystretch}{2em}
\multlinegap0pt
\usepackage{bm}
\usepackage{bbm}
\usepackage{dsfont}
\usepackage{xspace}
\usepackage{cancel}
\usepackage{mathtools}
\usepackage[shortlabels]{enumitem}
\usepackage{amsthm}
\makeatletter
\@ifundefined{lemma}{\newtheorem{lemma}{Lemma}}{}
\@ifundefined{theorem}{\newtheorem{theorem}[lemma]{Theorem}}{}
\makeatother
\newcommand{\R}{{\mathbb R}}
\newcommand{\cE}{{\mathcal E}}
\newcommand{\cR}{{\mathcal R}}
\newcommand{\vr}{{\varrho}}
\title[Nonlinear resistance forms]{Nonlinear resistance forms}
\author{Simon Puchert, Marcel Schmidt}
\date{Source: arXiv:2507.03426v1 (CC BY 4.0)}
\begin{document}
\maketitle

\noindent\emph{Source and licence.} Nonlinear resistance forms. Version arXiv:2507.03426v1 (CC BY 4.0), distributed under the Creative Commons Attribution 4.0 International licence, as recorded in the arXiv licence metadata for this exact version and shown on the abstract page https://arxiv.org/abs/2507.03426v1. Full rights notice: licenses/arxiv-2507.03426-CC-BY-4.0.txt.\par\smallskip

\noindent\textit{Editorial scope.} Counted unit: the Theorem whose source label is \texttt{theorem:triangle inequality} in the article (source0.tex lines 1781--1789, source label \texttt{theorem:triangle inequality}), delivered together with the article's own complete original proof of it (source0.tex lines 1790--1851, one \texttt{proof} environment of 62 lines). No mathematics is written, repaired or re-derived: statement and proof are the article's own text line for line. Required context retained uncounted: lemma:duality delta and nabla (lines 507--514). Every definition, display and label of those lines is retained unchanged. Omitted from this package and not redistributed: 1--506, 515--1780, 1852--2260. Nothing inside a retained excerpt is omitted or altered: every retained body is delivered exactly as the article's lines are written, so no omission receipt of any kind accompanies this package. Citations: the retained excerpts carry no citation command at all, so no bibliography entry of the article is transcribed and no reference is redistributed. Reference adaptation: none was needed and none was made; every reference inside the delivered unit resolves inside it, so no unresolved reference or citation is produced. Printed slips kept literal: no printed word, symbol, hypothesis or number of the counted statement or of its proof was altered. Layout and class: the article's own class is \texttt{amsart} with packages including amsmath, amssymb, amsfonts, amsthm, latexsym, textcomp, dsfont, babel, hyperref, cleveref, enumerate, inputenc, faktor, color; the wrapper is the lane's pinned \texttt{amsart} editorial wrapper at 10pt on A4, which restates the theorem-like environments the retained text needs and adds the article's own macro definitions that the retained text needs (\texttt{\textbackslash R}, \texttt{\textbackslash cE}, \texttt{\textbackslash cR}, \texttt{\textbackslash vr}), with extra wrapper packages (bm, bbm, dsfont, xspace, cancel, mathtools, enumitem). The delivered numbering is the unit's own. Assets: no retained span contains a figure environment, a graphics-inclusion command, a PDF-page inclusion, a table or a TikZ picture; no image file is copied into this package, the package declares no asset dependency and no image file is redistributed with it. Licence: version arXiv:2507.03426v1 of this article is distributed under the Creative Commons Attribution 4.0 International licence, as recorded in the arXiv licence metadata for this exact version and shown on the abstract page https://arxiv.org/abs/2507.03426v1. A full rights notice is delivered at licenses/arxiv-2507.03426-CC-BY-4.0.txt. Characters: every retained excerpt is pure ASCII, so the delivered engine, XeLaTeX with its default Latin Modern text font, reports no missing character.
 \begin{lemma}[Duality of the $\Delta_2$-condition and the $\nabla_2$-condition]\label{lemma:duality delta and nabla}
 \begin{enumerate}[(a)]
  \item   If $\vr$ satisfies the $\Delta_2$-condition, then $\vr^*$ satisfies the $\nabla_2$-condition.  
  \item If $\vr$ satisfies the $\nabla_2$-condition, then $\vr^*$ satisfies the $\Delta_2$-condition. 
  \item Assume that $\vr$ is lower semicontinuous. Then  $\vr$ satisfies the $\Delta_2$-condition if and only if $\vr^*$ satisfies the $\nabla_2$-condition and $\vr$ satisfies the $\nabla_2$-condition if and only if $\vr^*$ satisfies the $\Delta_2$-condition. 
 \end{enumerate}

 \end{lemma}

\begin{theorem}[Triangle inequality for the resistance]\label{theorem:triangle inequality}
Let $\cE$ be a nonlinear resistance form.  
\begin{enumerate}[(a)]
 \item The triangle inequality $\cR_t(x,z) \leq \cR_t(x,y) + \cR_t(y,z)$ holds for all $x,y,z \in X$ and $t > 0$. 
 \item If $\cE$ is symmetric, then  $\cR_t(x,y) = \cR_t(y,x)$ for all $x,y \in X$ and $t > 0$.
 \item If $\cE$ satisfies the $\nabla_2$-condition, then $\cR_t(x,y) < \infty$ for all $x,y \in X$ and $t> 0$.
\end{enumerate}
In particular, if $\cE$ is symmetric and satisfies the $\nabla_2$-condition, then for each $t > 0$ the function $\cR_t$ is a pseudometric. 
\end{theorem}

\begin{proof}
 (a):  We need to show
 $$ T: = t(f(x) - f(z)) - \mathcal{E}(f) \leq \mathcal{R}_t(x,y) + \mathcal{R}_t(y,z) $$
 %
 for all $f \in D(\cE)$. We have the trivial identity
 %
  \begin{align*}
 t(f(x) - f(z)) - \mathcal{E}(f)  & = t(f(x) - f(y)) + t(f(y) -  f(z)) - \mathcal{E}(f).
 \end{align*}
 %
 Hence, if $f(y) > f(x)$, then $T \leq \cR_t(y,z)$ and if $f(y) < f(z)$, then $T \leq \cR_t(x,y)$.
 
 It remains to treat the case $f(z) \leq f(y) \leq f(x)$.  Let $\alpha = f(y)$ and consider the normal contraction 
 %
   $$ C\colon \R\to\R, \quad C(t) = |t - \tfrac{\alpha}{2}| - |\tfrac{\alpha}{2}|. $$
 %
 Using the compatibility of $\cE$ with normal contractions, we obtain
 %
 \begin{align*}
  \mathcal{E}(f) & = \mathcal{E}(f) + \mathcal{E}(0) \\
                 & = \mathcal{E}(\tfrac{f}{2} + \tfrac{f}{2}) + \mathcal{E}(\tfrac{f}{2} - \tfrac{f}{2}) \\
                 & \geq \mathcal{E}(\tfrac{f}{2} + C(\tfrac{f}{2})) + \mathcal{E}(\tfrac{f}{2} - C(\tfrac{f}{2})) \\
                 & = \mathcal{E}(f\vee \alpha - \alpha_+) + \mathcal{E}(f\wedge \alpha + \alpha_-).
 \end{align*}
 %
 Hence, we can estimate 
 %
 \begin{align*}
  T & \leq  t(f(x) - f(y)) - \mathcal{E}(f\vee \alpha - \alpha_+)   + t(f(y) - f(z)) - \mathcal{E}(f\wedge \alpha + \alpha_-) \\
    & \leq \mathcal{R}_t(x,y) + \mathcal{R}_t(y,z).
 \end{align*}
%
The last inequality follows from the fact that $g = f \vee \alpha - \alpha_+$ and $h = f\wedge \alpha + \alpha_-$  satisfy $f(x) - f(y) = g(x) - g(y)$ and $f(y) - f(z) = h(y) - h(z)$ because $f(z) \leq \alpha = f(y) \leq f(x)$. 
%

(b): This is trivial.

(c): By  Lemma~\ref{lemma:duality delta and nabla} the $\nabla_2$-condition is equivalent to the $\Delta_2$-condition for $\cE^*$. In particular, $D(\cE^*)$ is a positive cone and we obtain $t (\delta_x - \delta_y) \in D(\cE^*)$ for all $t > 0$. 
%  
%  
%  
%  
%  We can assume without loss of generality that $g(x) \geq g(z)$ (else the inequality is trivial). By applying a suitable contraction $\phi\colon \R\to\R$, we can find a function $f = \phi g$ with $\mathcal{E}(f) \leq \mathcal{E}(g)$ and $f(x) - f(z) = g(x) - g(z)$ such that $f(x) \geq f \geq f(z)$. More specifically, we have
%   $$ \phi(t) = (g(z) \vee t \wedge g(x)) - \begin{cases} g(z) & 0 < g(z) \\ 0 & g(z) \leq 0 \leq g(x) \\ g(x) & g(x) < 0, \end{cases} $$
%  achieving the desired clamping of values.
%  
%  Now we have the situation $f(x) \geq c := f(y) \geq f(z)$. The next step splits the function at $c$ into the part that is cut off by $c$ and the remainder.
%  
%  Using the compatibility with the normal contraction
%   $$ C\colon \R\to\R, \quad C(x) = |x - \tfrac{c}{2}| - |\tfrac{c}{2}|, $$
%  we can show
%  
%  Each of these functions yields a lower bound for the resistances between $(x, y)$ and $(y, z)$, respectively. In total, we get
%  \begin{align*}
%   t\cdot(g(x) - g(z)) - \mathcal{E}(g) & \leq t\cdot(f(x) - f(z)) - \mathcal{E}(f) \\
%                                        & \leq t\cdot(f(x) - f(z)) - \mathcal{E}(f\vee c - c_+) - \mathcal{E}(f\wedge c + c_-) \\
%                                        & \leq t\cdot(f(x) - c) - \mathcal{E}(f\vee c - c_+) \\
%                                        & + t\cdot(c - f(z)) - \mathcal{E}(f\wedge c + c_-) \\
%                                        & \leq \mathcal{R}_t(x,y) + \mathcal{R}_t(y,z) \\
%  \end{align*}
%  Taking the supremum over all $g$ finishes the proof.
\end{proof}

\end{document}
